\section{Introduction}
\label{sec:introduction}

Complex reasoning problems often admit several solution paths, and a single
solver execution can fail even when another would succeed. Multi-sample
inference reduces this risk by generating several candidate solutions and
selecting or aggregating them, as in repeated sampling, self-consistency, and
verifier-guided
reasoning~\citep{chen2021codex,wang2023selfconsistency,cobbe2021verifiers}.
With a limited inference budget, the candidate pool succeeds when at least one
candidate solves the problem. Repeated sampling can substantially expand
problem coverage, although its gains depend on how candidates are generated
and selected~\citep{brown2024monkeys,snell2024scaling,wu2025inferencescaling}.
Recent training methods therefore score groups of outputs jointly or tailor
fine-tuning to multi-sample
selection~\citep{pmlr-v267-tang25o,walder2025pkpo,chow2025inferenceaware,chen2025rethinking}.
This line of work treats candidate-pool success as a training target rather
than only an inference-time metric.

A posterior over solver parameters naturally supplies the multiple candidates
used at inference. Each posterior draw instantiates a complete solver that
produces one candidate solution, so repeated draws assemble the inference pool
from a single learned distribution. We learn this distribution as a
regularized posterior, combining task loss with regularization toward a
prior~\citep{bissiri2016general,blundell2015weight,shen2024variational}.
The posterior's probability mass determines which solver variants appear in
the pool and how often they are sampled, so optimizing it directly shapes the
candidate generator used at inference.
Variational posterior learning typically optimizes the expected loss of one
sampled model. For multi-sample inference, we instead introduce finite-$K$
posterior learning, which trains the regularized posterior to make a finite
candidate pool more likely to contain at least one successful solution. Direct
optimization of this pool-level objective requires several posterior-sampled
solvers in every update, so its cost grows with the target pool size.

We reduce this cost with a Fenchel reformulation and a Bregman update that
require only one posterior draw per training update. The Fenchel reformulation
expresses the pool-level objective so that one sampled solver can update the
posterior when its failure is weighted by how likely the other candidates are
to fail. The required weight differs across problems and changes as the
posterior learns because sampled solvers do not fail equally often on every
problem. Estimating it from several fresh draws would cancel the computational
saving. We therefore maintain, for each training problem, a running summary of
how often its sampled solvers fail. A closed-form Bregman update refreshes this
summary after each observed failure, so the current summary guides the
posterior update and the new observation informs the next visit. The procedure
needs neither an inner optimization loop nor another solver execution. Each
update uses one sampled solver regardless of the candidate-pool size.

Across challenging reasoning tasks, our method improves candidate coverage
and, in a controlled comparison, achieves performance comparable to direct
multi-draw training with lower training cost. Our contributions are:

\begin{itemize}
  \item We formulate finite-$K$ posterior learning for reasoning, aligning a
  regularized posterior over solver parameters with the collective success of
  the finite candidate pool used at inference.

  \item We derive an exact Fenchel reformulation of the differentiable
  finite-$K$ objective that enables posterior updates from one sampled solver,
  guided by a problem-specific summary of posterior failures.

  \item We derive a closed-form Bregman update for tracking these failure
  summaries and analyze how sampling noise and posterior drift affect their
  tracking error and the resulting finite-$K$ objective gap.
\end{itemize}
